\documentclass[12pt]{article}

% Layout.
\usepackage[top=1in, bottom=0.75in, left=0.75in, right=0.75in, headheight=1in, headsep=6pt]{geometry}

% Fonts.
\usepackage{mathptmx}

\usepackage[scaled=0.86]{helvet}
\renewcommand{\emph}[1]{\textsf{\textbf{#1}}}

% Misc packages.
\usepackage{amsmath,amssymb,latexsym,amsthm}
\usepackage{graphicx,hyperref}
\usepackage{array}
\usepackage{xcolor}
\usepackage{multicol}
\usepackage{tabularx,colortbl}
\usepackage{enumitem}
\usepackage{soul,multicol}
\usepackage{setspace}


\hypersetup{
    colorlinks=true,
    linkcolor=blue,
    filecolor=magenta,      
    urlcolor=blue,
    pdftitle={Proofs Worksheet},
    pdfpagemode=FullScreen,
    }

\def\mailto#1{\href{mailto:#1}{#1}}

%% special math symbols
\newcommand{\N}{\mathbb{N}}
\newcommand{\Z}{\mathbb{Z}}
\newcommand{\R}{\mathbb{R}}
\newcommand{\Q}{\mathbb{Q}}
\newcommand{\sse}{\subseteq}
\newcommand{\mcp}{\mathcal{P}}
\newcommand{\ol}[1]{\overline{#1}}

%other specials
\newcommand{\be}{\begin{enumerate}}
\newcommand{\ee}{\end{enumerate}}
\newcommand{\se}{\subseteq}
\newcommand{\spe}{\supseteq}
\newcommand{\ds}{\displaystyle}
\newcommand{\lcm}{\text{lcm}}
%\newcommand{\gcd}{\text{gcd}}

% Paragraph spacing
\parindent 0pt
\parskip 6pt plus 1pt
\def\tableindent{\hskip 0.5 in}
\def\ts{\hskip 1.5 em}

%header
\usepackage{fancyhdr}
\pagestyle{fancy} 
\lhead{\large\sf\textbf{MATH 265: Introduction to Mathematical Proofs}}
\rhead{\large\sf\textbf{Worksheet 22: Ch 12}}

\newcommand{\localhead}[1]{\par\smallskip\textbf{#1}\nobreak\\}%
\def\heading#1{\localhead{\large\emph{#1}}}
\def\subheading#1{\localhead{\emph{#1}}}

\newenvironment{clist}%
{\bgroup\parskip 0pt\begin{list}{$\bullet$}{\partopsep 4pt\topsep 0pt\itemsep -2pt}}%
{\end{list}\egroup}%

\begin{document}
\begin{enumerate}
\item State each definition and draw helpful pictures.
	\begin{enumerate}
	\item The relation $R$ from $A$ to $B$ is a function if\\
	\item The function $f$ from $A$ to $B$ injective if \\
	\item The function $f$ from $A$ to $B$ surjective if \\
	\end{enumerate}
\item Let $f: \R-\{0\} \to \R -\{1\}$ be a function defined by $f(x)=\frac{x+1}{x}.$ Prove that $f$ is injective and surjective.
\vfill
\item Let $f:\Z \times \Z \to \Z$ be a function defined by $f(m,n)=2m-6n.$ Prove that $f$ is neither injective nor surjective.
\vfill
\newpage
\item For each function below, determine if it is injective or surjective and prove your answer is correct. If the function is not surjective, determine its range.
	\begin{enumerate}
	\item  $f:\Z \times \Z \to \Z$ defined as $f(m,n)=2m-3n.$
	%not 1-1, onto
	\vfill
	\item $\theta: \{0,1\} \times \N \to \Z$ defined as $\theta(a,b)=(-1)^ab.$
	%not onto, is 1-1
	\vfill
	\item $g: \N \to \Z$ defined as $g(n)=\begin{cases} \frac{n}{2} & n \text{ even} \\
	(-1)\frac{n-1}{2} & n \text{ odd} \\
	\end{cases}$
	
	\vfill
	\end{enumerate} 
\end{enumerate}
\end{document}










